\section{A twelve-variable companion of degree 629}
\label{sec:twelve}

The same integrality mechanism also improves the degree of a construction
with fewer supplied variables. The result in this section uses twelve strictly
positive integers and has exact degree $629$, with $114$ live binary arithmetic
operations. Its three-equation version has maximum residual degree $622$ and
$109$ operations, or $108$ after the established positive transport shift.
The comparison is with package~52's twelve-positive lean equation, of degree
$897$ and $108$ operations. Thus the new point reduces degree by approximately
$29.9\%$, at a cost of six operations. It does not reduce the number of supplied
variables, and no global optimality or literature-wide novelty is claimed.

All notation in this section is local. The four fixed ports are the same
complete authentic tuple $(\beta,K_0,M_C,M_F)$ used elsewhere in this report;
$M_F$ is already shifted. In particular $\beta=2^{590069}-1$. Canonical
completeness uses odd native exponent $N$, so its old coordinate satisfies
$J_{\rm old}\equiv1\pmod3$. We may therefore put $J_{\rm old}=3J+1$,
retaining completeness with positive $J$ by taking $N>1$. This changes the
parameterization, not the coefficient tuple. It gives $q\equiv2\pmod3$.
The point of this residue restriction is that the later divisor $q$ is
coprime to $3$.

\subsection{The fixed polynomial}

The twelve supplied coordinates are
\[
 J,F,\alpha,z,h,i,T,s,w,\eta,\zeta,t\in\mathbb Z_{>0}.
\]
Define the following polynomial abbreviations:
\begin{align*}
 j&=3J+1,&q&=\beta j+1,&C&=q-F-\alpha,\\
 X&=wq,&Y&=sq,&E&=XY,\\
 k&=\eta+\zeta,&c&=kY+\eta,&a&=E+Y,\\
 H&=4a+3,&\Delta&=a^2+H=(a+2)^2-1,\\
 \ell&=EkY,&A&=1+\ell(\ell+k),&G&=1+\Delta c^2,\\
 b_0&=X+ac,&&&
\end{align*}
\[
 R=\bigl((q-1)(q-F)+\alpha\bigr)(q^2-1)+(M_C+qM_F)j.
\]
\begin{align*}
 N_k&=k-hE-R,\\
 N_t&=(K_0+w)C+q-F-z(q-1).
\end{align*}
\begin{align*}
 I&=ic^2,&D&=1+\Delta I^2,&Q&=(\Delta I)^2,\\
 d_0&=Q-1,&v&=cT,&b&=c+RD,\\
 e&=Q(Dv^2+b^2)-1,&B_a&=2Qvb.
\end{align*}
The new scaled carrier is
\begin{align*}
 S&=q^2H,&L&=q^2(1+HIT),\\
 W&=q^2b_0+H\bigl[q^2(t+b+\ell)+s\bigr].
\end{align*}
Let
\begin{equation}
 \mathcal N=
 \prod_{\substack{g=\pm\sqrt G,\ \tau=\pm\sqrt A\\f=\pm\sqrt D}}
 \left[d_0(W-Lg-S\tau)^2-S^2e+S^2B_af\right].
 \label{eq:twelve-norm}
\end{equation}
Every formal sign is included. Consequently this is an integer polynomial;
no radicals or intermediate registers are additional unknowns. Formal signs
are not being identified with independent field automorphisms after
specialization. The equation is
\begin{equation}
 \boxed{P_{12}=N_kN_t(1+\mathcal N)-1=0.}
 \label{eq:twelve-final}
\end{equation}

For an explicit division-free evaluator, put
\begin{align*}
 d_A&=d_0S^2A,&U_Q&=S^2Q,\\
 P_H&=U_QDv^2,&U_H&=U_Qb^2,\\
 e_H&=P_H+U_H-S^2,&j_H&=4P_HU_H,\\
 \mathcal T&=2d_0(W^2+L^2G),&
 \kappa&=[d_0(W^2-L^2G)]^2,\\
 \mathcal A&=2(d_A+e_H),&
 \mathcal B&=(d_A-e_H)^2-j_H,\\
 a_r&=\mathcal T-\mathcal A,&b_r&=\mathcal B-\kappa,\\
 v_r&=16d_Aj_H-a_r(a_r\mathcal T+2b_r),&
 u_r&=b_r^2-a_r^2\kappa,\\
 \mathcal N&=\kappa v_r^2-u_r(\mathcal T v_r-u_r).
\end{align*}
This is the inherited quadratic-remainder identity for the formal norm.
The supplied JSON circuit gives every live gate and includes all
multiplications by fixed numerals. The system version compares the final
norm terms directly and imposes $N_k=N_t=1$.

\subsection{Bounds before recovering divisibility}

The predecessor norm identity gives
$\mathcal N\equiv(u_r-x_0v_r)^2\pmod4$, where
$x_0=d_0(W^2+L^2G)$. At a root of \eqref{eq:twelve-final}, all three factors
are integer units. The option $\mathcal N=-2$ is impossible modulo~$4$.
Thus $\mathcal N=0$ and $N_k=N_t=\varepsilon\in\{1,-1\}$.

The authentic native-mask inequalities applied to the transport unit give
\[
 C\ge0,\qquad q-F=C+\alpha>0,\qquad q\ge16,
 \qquad 3q<R<q^4-2.
\]
The boundary $C=0$ is not discarded at this stage. We cannot yet use
$q^3\mid Y$ or $E>R$: the present definition is only $Y=sq$. Instead,
\begin{equation}
 X,Y\ge q,\quad E\ge q^2,\quad
 k=hE+R+\varepsilon>R,\quad c>RY,\quad
 a+2<2E,\quad I\ge c^2.
 \label{eq:twelve-early}
\end{equation}
In bounds below, $g,f,\tau$ denote positive real square roots. Then
$g<2Ec$, $f<2EI$, and the inherited elementary expansion gives
\[
 \Delta c^2-A>
 k^2\bigl[EY(2Y^2-1)+Y^4\bigr]-1>0.
\]
In particular
\begin{equation}
 1<A<\Delta c^2\le\Delta I=\sqrt Q.
 \label{eq:twelve-exception-bound}
\end{equation}

The earlier lean construction used $E>R$ in its outer-root estimate. Here it
can be replaced by a consequence of $c>RY$. Since
$D<4E^2I^2$, $c<I^2$, and $d_0>\tfrac12E^4I^2$, we have
\[
 b+cf<(4R+1)E^2I^2,
 \qquad \frac{b+cf}{d_0}<\frac{9R}{E^2}.
\]
On the other hand, the exact identity
$(Ig)^2-(cf)^2=I^2-c^2$ implies
\[
 Ig-cf=\frac{I^2-c^2}{Ig+cf}>\frac{I}{8Ec}
       =\frac{ic}{8E}.
\]
Finally $icE>RYE=RXY^2>72R$ by \eqref{eq:twelve-early}. Therefore
\begin{equation}
 \frac{b+cf}{d_0}<\frac{9R}{E^2}<\frac{ic}{8E}<Ig-cf.
 \label{eq:twelve-new-bound}
\end{equation}
This is the extra estimate that makes divisibility recovery possible without
assuming the scale one is trying to recover.

\subsection{Integral roots force the missing scale}

For a coherent auxiliary root $y$ with $(f')^2=D$, the norm relation is
\[
 d_0y^2=Q(vf'-b)^2-1.
\]
The standard square estimate, also valid for complex $y$, gives
\[
 |y|<b+Tcf+T\frac{b+cf}{d_0}<b+ITg.
\]
Consequently any root $x=\tau'+y$ of the inner polynomial has
$|x|<\tau+b+ITg$. But the negative-main-root branch would give
\[
 \frac{W+Lg}{S}
 =b+ITg+\left(t+\frac{s}{q^2}\right)
   +\ell+\frac{b_0+g}{H}.
\]
This exceeds that upper bound: $t+s/q^2>0$,
$\tau<\ell+k$, and $(b_0+g)/H>c/3>k$, since $Y\ge16$.
Actual quadratic conjugation would send an irrational positive-main-root
branch to this forbidden branch. Hence $g$ is a positive integer and
$x=(W-Lg)/S$ is rational.

For completeness, the inherited inner-root reconstruction can be stated
explicitly. Define
\begin{align*}
 C_x&=d_0(x^2+A)-e,\\
 F(x)&=\bigl[C_x^2+DB_a^2-4d_0^2x^2A\bigr]^2-4C_x^2DB_a^2.
\end{align*}
The norm implies $F(x)=0$. If $C_x=0$, eliminating $x^2$ would give
\[
 d_0^2A^2-d_0eA+Q^2v^2b^2D=0,
\]
which forces $Q\mid A(A-1)$, contrary to
\eqref{eq:twelve-exception-bound}. Likewise $F(0)=0$ would force
$Q\mid(A-1)^2$, also impossible. Thus the denominators in
\begin{align*}
 U_x&=C_x^2+DB_a^2-4d_0^2x^2A,\\
 f'&=-\frac{U_x}{2C_xB_a},&
 \tau'&=\frac{C_x+B_af'}{2d_0x},&y&=x-\tau'
\end{align*}
are nonzero. The identity $F(x)=0$ gives
$(f')^2=D$, $(\tau')^2=A$, and the auxiliary Pell equation above.
The rational square roots $f',\tau'$ are integers. The inherited rational
Pell saturation lemma says that, for integers $m\ge2,V$ and rational $y$,
\[
 m^2V^2-(m^2-1)y^2=1\quad\Longrightarrow\quad y\in\mathbb Z.
\]
Applying it with $m=\Delta I$ and $V=vf'-b$ shows that $y$, and hence $x$,
is integral. The lemma's ordinary proof is preserved in the upstream proof
files; it is not a newly supplied proof-assistant certificate.

Reduce $W=Lg+Sx$ modulo $q^2$. This gives $q^2\mid Hs$.
But $Y=sq$ gives $H\equiv3\pmod q$, while $q\equiv2\pmod3$.
Thus $\gcd(q,H)=1$ and
\begin{equation}
 q^2\mid s,
 \qquad s_{\rm old}=s/q^2\in\mathbb Z_{>0},
 \qquad Y=s_{\rm old}q^3.
 \label{eq:twelve-divisibility}
\end{equation}

There is now a short soundness proof using the old equation directly.
Let $\mathcal N_{\rm lean}$ be package~52's unscaled twelve-variable lean norm,
on the same fixed tuple. In the localization at $q$ there is the exact identity
\begin{equation}
 \mathcal N=q^{32}
   \mathcal N_{\rm lean}\left(t+\frac{s}{q^2}\right),
 \qquad Y=(s/q^2)q^3.
 \label{eq:twelve-identity}
\end{equation}
Each of eight factors gains $q^4$. Its carrier degree is at most $16$, so
all denominators clear. By \eqref{eq:twelve-divisibility}, the mapped old
coordinates $s_{\rm old}$ and $t_{\rm old}=t+s_{\rm old}$ are positive
integers. The index and transport factors are unchanged. Thus a new root
produces a root of the established old positive equation. Its compiler and
K215 calibration recover a finite ZFC refutation.

\subsection{Completeness and coefficient-preserving equivalence}

Choose canonical normalized compiler witnesses with odd $N>1$. Keep the
retained coordinates, put $J=(J_{\rm old}-1)/3$, $s=s_{\rm old}q^2$, and set
\[
 t=\gamma+ITg+\tau+y-b-\ell-s_{\rm old}.
\]
All coordinates are integral. The auxiliary Pell equation gives
$y\ge b-vf$, and $Ig>cf$. Also
\[
 \tau-\ell=\frac{\ell k+1}{\tau+\ell}>\frac{k}{3},
 \qquad
 \frac{k}{3}>\frac{XY}{3}>\frac{Y}{q^3}=s_{\rm old}.
\]
Here $\ell>k$, $\tau<\ell+k$, $k>XY$, and $X\ge q\ge16$ are canonical
compiler inequalities. Therefore
\[
 t\ge\gamma+T(Ig-cf)+(\tau-\ell)-s_{\rm old}>0.
\]
A coherent factor in \eqref{eq:twelve-norm} vanishes, and $N_k=N_t=1$.
This proves both existence implications. Composing with the inherited
syntactic-refutation calibration yields, by ordinary mathematics formalizable
in ZFC,
\[
 \mathrm{ZFC}\vdash
 \bigl[\text{$P_{12}$ has no positive integer zero}\bigr]
 \longleftrightarrow\mathrm{Con}(\mathrm{ZFC}).
\]
No truth value of either equivalent assertion is proved. The separately
supplied system replaces the product by $N_k=N_t=1$ and $\mathcal N=0$.
Its shifted transport version uses $(K_0+w)C=F+z(q-1)$; the established
positive transport shift preserves existence after the scale has been
recovered.

\subsection{Exact degree and reproducible evidence}

Set the leading homogeneous forms
\begin{align*}
 q_0&=3\beta J,&X_0&=wq_0,&Y_0&=sq_0,&E_0&=X_0Y_0,\\
 k_0&=\eta+\zeta,&c_0&=k_0Y_0,&I_0&=ic_0^2,&\ell_0&=E_0c_0,\\
 H_0&=4E_0,&Q_0&=E_0^4I_0^2,&R_0&=q_0^3(q_0-F),\\
 B_0&=R_0E_0^2I_0^2,&V_0&=I_0T\ell_0.
\end{align*}
The respective degrees of $H_0,\ell_0,Q_0,B_0,V_0,c_0$ are
$4,7,30,26,15,3$. The exact leading cancellations in the unscaled lean
norm give the following eight factor degrees:
\[
 79,79,79,79,71,71,67,65.
\]
Their sum is $590$. Four opposite $f/g$ branches have leading form
$-\operatorname{sgn}(g)4H_0^2Q_0B_0V_0$; the two matched branches with
negative $\tau$ have leading form $4H_0^2Q_0B_0\ell_0$; the other two
matched branches have leading forms
\[
 H_0^2Q_0B_0c_0,
 \qquad
 H_0^2Q_0B_0\left(2t-k_0-\frac{iTc_0}{E_0}\right).
\]
Their product is
\begin{equation}
 2^{12}H_0^{16}Q_0^8B_0^8V_0^4\ell_0^2c_0
 \left(2t-k_0-\frac{iTc_0}{E_0}\right).
 \label{eq:twelve-leading}
\end{equation}
The displayed denominator cancels against preceding factors. The result is
a nonzero polynomial, since $t$ is independent. These are the same exact
lean-norm cancellations as in the bundled predecessor proof, now with
$\deg X=\deg Y=2$; in particular the cancellation uses
$(Ig)^2-(cf)^2=I^2-c^2$, not independent treatment of these square roots.

The shift $s/q^2$ in \eqref{eq:twelve-identity} has degree $-1$, below
$\deg t=1$, so it does not change \eqref{eq:twelve-leading}. The factor
$q^{32}$ adds $32$ degrees, giving $\deg\mathcal N=622$. The index factor
has nonzero degree-$5$ leader $-hE_0$, and the transport factor has nonzero
degree-$2$ leader $w(q_0-F-\alpha)-zq_0$. Hence
\[
 \deg P_{12}=622+5+2=629.
\]
This is uniform for every authentic positive $\beta$ satisfying the stated
coefficient and completeness interface.

The companion checker performs $60$ modular formal-norm multiplications,
$59$ valid modular checks of \eqref{eq:twelve-identity}, and six complete
univariate polynomial computations: three with formal-port fixtures and
three with the actual fixed tuple. All six give degrees $622$ and $629$.
It also checks all four compressed and uncompressed coefficient hashes and
all fourteen bundled source-file pins. These are finite algebra diagnostics;
the unbounded equivalence rests on the preceding argument and the inherited
compiler theorems. The files \texttt{twelve/upstream/LEAN\_PROOF.md},
\texttt{INNER\_RECOVERY\_PROOF.md}, \texttt{ELEVEN\_POSITIVE\_PROOF.md}, and
\texttt{ORDINARY\_BLANK\_PROOF.md} preserve those ordinary mathematical
boundaries explicitly.
